\paragraph{Another statement can reject a cash-compatible history.}
The preceding borrowing examples leave the other disclosed quantities
unconstrained. That omission matters: mapping their assumed capital balance
to the reported closing equity rejects the annual candidate, although the
cash-history system and the balance-statement system are each feasible by
themselves. Compatibility must be checked on their intersection, using one
shared history, rather than inferred from separate successful calculations.

The optional \texttt{RetainedStatementHistoryAttachment/1} adds selected
balance equations to the same history-and-cash constraints. It preserves
the original source coefficients and contexts, moves the snapshot date into
the explicit time type, and requires an explicit signed field mapping.
All selected equations remain even when some operands have no history
mapping; those unmapped quantities create unresolved obligations. The
complete source report and both parent assessments remain attached. A
different basis, monetary scale or quantity cannot be substituted because
two amounts happen to agree.

In each annual and interim example, 27 reported closing-balance quantities
join the three cash observations through ten selected relations. A
constructed compatible history uses one signed account for each disclosed
leaf caption and explicit sums for totals. It assumes unchanged noncash
balances and specifies opening capital by reversing a hypothetical net
capital contribution or return from the reported closing caption. These
assumptions explain the construction; they are not recovered opening facts.
Two gross contribution-and-return sequences can again fit the same selected
closing values. Fourteen complete exchanges, 13 focused tests and 62 existing
tests check the joint constraints, conflicting and unresolved cases, and
exact replay. The complete exchanges are retained in
\begin{quote}\small\ttfamily
data/accounting-master/\\
retained-statement-history-20261025/
\end{quote}
The reported opening balance sheet and equity movements provide further
constraints that these examples have not yet used. They can reject the
assumed opening or event sequence even after every selected closing caption
agrees. The remaining task is therefore to attach those reported movements,
while continuing to distinguish a conditional witness from actual history.
